\setchapterpreamble[u]{\margintoc}
\chapter{Sample Spaces, Conditioning, and Bayes}
\labch{sample-spaces}

\sources{Bertsekas and Tsitsiklis, \emph{Introduction to Probability}, 2nd ed.\
\cite{bertsekas2008probability}, Chapter~1; MIT 6.041 \cite{ocw6041},
Lectures~1--4; MIT 18.05 \cite{ocw1805}.}

Part~III covers the probability that a learning problem needs: enough to state
what a model assumes about its data, and enough to derive the two loss functions
that dominate practice. It stops at estimation, as agreed.

\section{Sample spaces and the axioms}
\labsec{axioms}

An experiment is described by a \emph{sample space} \(\Omega\), the set of
possible outcomes. An \emph{event} is a subset of \(\Omega\). A probability
measure assigns to each event a number \(\Pr(A)\) subject to three requirements:

\begin{enumerate}
	\item \(\Pr(A)\ge0\);
	\item \(\Pr(\Omega)=1\);
	\item \(\Pr\bigl(\bigcup_i A_i\bigr)=\sum_i\Pr(A_i)\) for disjoint events.
\end{enumerate}

\marginnote{Everything else is bookkeeping on top of these three lines ---
including the results in \refch{limit-theorems} that justify learning from
samples at all.}

Immediate consequences are \(\Pr(A^{c})=1-\Pr(A)\), monotonicity, and the
inclusion--exclusion formula
\(\Pr(A\cup B)=\Pr(A)+\Pr(B)-\Pr(A\cap B)\).

\section{Conditional probability}
\labsec{conditional-probability}

\begin{definition}
For \(\Pr(B)>0\), the probability of \(A\) given \(B\) is
\[
	\Pr(A\mid B)=\frac{\Pr(A\cap B)}{\Pr(B)}.
\]
\end{definition}

Conditioning is renormalisation: \(B\) becomes the new sample space, and the
measure is rescaled so that it again totals one. For fixed \(B\) the map
\(A\mapsto\Pr(A\mid B)\) satisfies the three axioms, so every result proved for
probabilities holds for conditional probabilities as well.

Rearranging gives the multiplication rule
\(\Pr(A\cap B)=\Pr(A\mid B)\Pr(B)\), and chaining it,
\[
	\Pr(A_1\cap\cdots\cap A_n)
	=\Pr(A_1)\Pr(A_2\mid A_1)\cdots\Pr(A_n\mid A_1\cap\cdots\cap A_{n-1}).
\]

\section{Total probability and Bayes' rule}
\labsec{bayes}

Let \(B_1,\ldots,B_k\) partition \(\Omega\). Splitting any event along the
partition gives the \emph{total probability} formula
\[
	\Pr(A)=\sum_{i=1}^{k}\Pr(A\mid B_i)\Pr(B_i),
\]
and dividing the multiplication rule by it gives \emph{Bayes' rule},
\begin{equation}
	\Pr(B_j\mid A)=\frac{\Pr(A\mid B_j)\Pr(B_j)}{\sum_{i}\Pr(A\mid B_i)\Pr(B_i)} .
	\labeq{bayes}
\end{equation}

\begin{example}
A test for a condition present in \(0.1\%\) of a population detects it in
\(99\%\) of carriers and returns a false positive for \(1\%\) of non-carriers.
For a positive result,
\[
	\Pr(\text{carrier}\mid+)
	=\frac{0.99\times0.001}{0.99\times0.001+0.01\times0.999}
	=\frac{0.00099}{0.01098}\approx0.090 .
\]
A test that is right \(99\%\) of the time leaves the subject about \(9\%\) likely
to be a carrier, because the prior is small enough that false positives
outnumber true ones by roughly ten to one.
\end{example}

\marginnote{The same arithmetic explains why a classifier with excellent accuracy
on a rare class can still produce mostly wrong positive predictions. Accuracy and
precision answer different questions; \refeq{bayes} is the conversion between
them.}

\section{Independence}
\labsec{independence}

\begin{definition}
Events \(A\) and \(B\) are \emph{independent} when
\(\Pr(A\cap B)=\Pr(A)\Pr(B)\), equivalently \(\Pr(A\mid B)=\Pr(A)\) when
\(\Pr(B)>0\).
\end{definition}

Independence is an assumption about the mechanism generating the data, not a
property that can be read off a formula. It is also fragile: events may be
pairwise independent without being jointly independent, and independence
generally fails once one conditions on a third event.

\begin{remark}
The i.i.d.\ assumption of \refch{learning-problem} is exactly independence
applied to the examples in a dataset, plus the requirement that they share one
distribution. Nearly every guarantee in \refch{limit-theorems} rests on it, and
distribution shift is the name for what happens when the second half fails.
\end{remark}
